% GENERATED FILE. Edit canonical lessons, then run npm run generate:books.
% canonical-source-hash: fnv1a64:5e2bc71064545be7
% canonical-lessons: GU-C437-thar, GU-C437-jhakal, GU-C437-tadh, GU-C437-thijvu, GU-C437-thartharvu

\chapter{Frost and the Cold}
\label{ch:gu-frost-and-the-cold}
\hlchaptermodality{437}

\begin{chapteropening}
\textbf{By the end of this chapter:} \emph{I can talk about a cold winter night.}

Complete the last lesson of chapter 437: I can talk about a cold winter night.
\end{chapteropening}

\section[\texorpdfstring{\gu{ઠાર} (ṭhār)}{ઠાર (ṭhār)}]{\gu{ઠાર} (ṭhār) — a hard frost; a biting chill}

\label{lesson:GU-C437-thar}



\begin{quote}
Before the new one: say the Gujarati for to scorch, then the Gujarati for to bake (in the heat).
\end{quote}

\subsection*{You'll want to know: \gu{ઠાર}}
\textbf{\gu{ઠાર}} — \emph{ṭhār} — \textquotedblleft{}a hard frost; a biting chill\textquotedblright{}.

\begin{quote}
\gu{ગઈકાલે} \gu{રાત્રે} \gu{ઘણો} \gu{ઠાર} \gu{પડ્યો}.
\end{quote}

\begin{quote}
\emph{gaīkāle rātre ghaṇo ṭhār paḍyo.}
\end{quote}

\begin{quote}
There was a hard frost last night.
\end{quote}

A masculine noun.

\begin{cousinweb}
\textbf{\gu{ઠાર}} (\emph{ṭhār}) goes back to \emph{*sthāra}, \textquotedblleft{}firm, firmness\textquotedblright{}, a form CDIAL reconstructs; CDIAL lists the Gujarati word, \textquotedblleft{}frost\textquotedblright{} (CDIAL 13765): the cold that makes things set hard.
\end{cousinweb}

\begin{grammarlens}[title={What you've built}]
One more word for getting around.
\end{grammarlens}

\subsection*{Your turn}
\begin{itemize}
\raggedright
  \item \emph{Say it:} \emph{ṭhār}
  \item \emph{Say it:} \emph{ṭhār}, once more
  \item \emph{Say it:} \emph{gaīkāle rātre ghaṇo ṭhār paḍyo.} — There was a hard frost last night.
  \item \emph{Recall:} say the Gujarati for to scorch, then the Gujarati for to bake (in the heat), then say \emph{ṭhār} again
\end{itemize}

\begin{tcolorbox}[breakable,colback=teal!4,colframe=teal!35!black,title={Before you move on}]
What does \gu{ઠાર} mean? (\textquotedblleft{}A hard frost; a biting chill\textquotedblright{}.) Say it once more.
\end{tcolorbox}
\section[\texorpdfstring{\gu{ઝાકળ} (jhākaḷ)}{ઝાકળ (jhākaḷ)}]{\gu{ઝાકળ} (jhākaḷ) — heavy dew; morning mist}

\label{lesson:GU-C437-jhakal}



\begin{quote}
Before the new one: say the Gujarati for to bake (in the heat), then the Gujarati for a biting chill.
\end{quote}

\subsection*{You'll want to know: \gu{ઝાકળ}}
\textbf{\gu{ઝાકળ}} — \emph{jhākaḷ} — \textquotedblleft{}heavy dew; morning mist\textquotedblright{}.

\begin{quote}
\gu{સવારે} \gu{ઘાસ} \gu{પર} \gu{ઝાકળ} \gu{હતું}.
\end{quote}

\begin{quote}
\emph{savāre ghās par jhākaḷ hatũ.}
\end{quote}

\begin{quote}
In the morning there was heavy dew on the grass.
\end{quote}

A neuter noun.

\begin{cousinweb}
\textbf{\gu{ઝાકળ}} (\emph{jhākaḷ}) No dictionary we use records an older form, so its history is not settled.
\end{cousinweb}

\begin{grammarlens}[title={What you've built}]
One more word for getting around.
\end{grammarlens}

\subsection*{Your turn}
\begin{itemize}
\raggedright
  \item \emph{Say it:} \emph{jhākaḷ}
  \item \emph{Say it:} \emph{jhākaḷ}, once more
  \item \emph{Say it:} \emph{savāre ghās par jhākaḷ hatũ.} — In the morning there was heavy dew on the grass.
  \item \emph{Recall:} say the Gujarati for to bake (in the heat), then the Gujarati for a biting chill, then say \emph{jhākaḷ} again
\end{itemize}

\begin{tcolorbox}[breakable,colback=teal!4,colframe=teal!35!black,title={Before you move on}]
What does \gu{ઝાકળ} mean? (\textquotedblleft{}Heavy dew; morning mist\textquotedblright{}.) Say it once more.
\end{tcolorbox}
\section[\texorpdfstring{\gu{ટાઢ} (ṭāḍh)}{ટાઢ (ṭāḍh)}]{\gu{ટાઢ} (ṭāḍh) — the cold; cold weather}

\label{lesson:GU-C437-tadh}



\begin{quote}
Before the new one: say the Gujarati for a biting chill, then the Gujarati for heavy dew.
\end{quote}

\subsection*{You'll want to know: \gu{ટાઢ}}
\textbf{\gu{ટાઢ}} — \emph{ṭāḍh} — \textquotedblleft{}the cold; cold weather\textquotedblright{}.

\begin{quote}
\gu{આજે} \gu{ઘણી} \gu{ટાઢ} \gu{છે}.
\end{quote}

\begin{quote}
\emph{āje ghaṇī ṭāḍh che.}
\end{quote}

\begin{quote}
It is very cold today.
\end{quote}

A feminine noun.

\begin{cousinweb}
\textbf{\gu{ટાઢ}} (\emph{ṭāḍh}) goes back to Sanskrit \emph{stabdhi}, \textquotedblleft{}fixedness, numbness\textquotedblright{}; CDIAL lists the Gujarati word under it (CDIAL 13677): the cold that makes your hands stiff.
\end{cousinweb}

\begin{grammarlens}[title={What you've built}]
One more word for getting around.
\end{grammarlens}

\subsection*{Your turn}
\begin{itemize}
\raggedright
  \item \emph{Say it:} \emph{ṭāḍh}
  \item \emph{Say it:} \emph{ṭāḍh}, once more
  \item \emph{Say it:} \emph{āje ghaṇī ṭāḍh che.} — It is very cold today.
  \item \emph{Recall:} say the Gujarati for a biting chill, then the Gujarati for heavy dew, then say \emph{ṭāḍh} again
\end{itemize}

\begin{tcolorbox}[breakable,colback=teal!4,colframe=teal!35!black,title={Before you move on}]
What does \gu{ટાઢ} mean? (\textquotedblleft{}The cold; cold weather\textquotedblright{}.) Say it once more.
\end{tcolorbox}
\section[\texorpdfstring{\gu{થીજવું} (thījvũ)}{થીજવું (thījvũ)}]{\gu{થીજવું} (thījvũ) — to congeal; to freeze solid}

\label{lesson:GU-C437-thijvu}



\begin{quote}
Before the new one: say the Gujarati for heavy dew, then the Gujarati for the cold.
\end{quote}

\subsection*{You'll want to know: \gu{થીજવું}}
\textbf{\gu{થીજવું}} — \emph{thījvũ} — \textquotedblleft{}to congeal; to freeze solid\textquotedblright{}.

\begin{quote}
\gu{રાત્રે} \gu{ઠારમાં} \gu{ઘી} \gu{થીજી} \gu{ગયું}.
\end{quote}

\begin{quote}
\emph{rātre ṭhārmā̃ ghī thījī gayũ.}
\end{quote}

\begin{quote}
In the night's frost the ghee set solid.
\end{quote}

Past (it …): \gu{તે} \gu{પાણી} \gu{થીજી} \gu{ગયું}, \emph{te pāṇī thījī gayũ}.

\begin{cousinweb}
\textbf{\gu{થીજવું}} (\emph{thījvũ}) goes back to \emph{*sthiryati}, \textquotedblleft{}becomes firm\textquotedblright{}, a form CDIAL reconstructs; CDIAL lists the Gujarati verb, \textquotedblleft{}to thicken from cold\textquotedblright{} (CDIAL 13772).
\end{cousinweb}

\begin{grammarlens}[title={What you've built}]
One more word for getting around.
\end{grammarlens}

\subsection*{Your turn}
\begin{itemize}
\raggedright
  \item \emph{Say it:} \emph{thījvũ}
  \item \emph{Say it:} \emph{thījvũ}, once more
  \item \emph{Say it:} \emph{rātre ṭhārmā̃ ghī thījī gayũ.} — In the night's frost the ghee set solid.
  \item \emph{Recall:} say the Gujarati for heavy dew, then the Gujarati for the cold, then say \emph{thījvũ} again
\end{itemize}

\begin{tcolorbox}[breakable,colback=teal!4,colframe=teal!35!black,title={Before you move on}]
What does \gu{થીજવું} mean? (\textquotedblleft{}To congeal; to freeze solid\textquotedblright{}.) Say it once more.
\end{tcolorbox}
\section[\texorpdfstring{\gu{થરથરવું} (thartharvũ)}{થરથરવું (thartharvũ)}]{\gu{થરથરવું} (thartharvũ) — to shiver; to tremble all over}

\label{lesson:GU-C437-thartharvu}



\begin{quote}
Before the new one: say the Gujarati for the cold, then the Gujarati for to congeal.
\end{quote}

\subsection*{You'll want to know: \gu{થરથરવું}}
\textbf{\gu{થરથરવું}} — \emph{thartharvũ} — \textquotedblleft{}to shiver; to tremble all over\textquotedblright{}.

\begin{quote}
\gu{ટાઢમાં} \gu{બાળકો} \gu{થરથરતાં} \gu{હતાં}.
\end{quote}

\begin{quote}
\emph{ṭāḍhmā̃ bāḷko tharthartā̃ hatā̃.}
\end{quote}

\begin{quote}
The children were shivering in the cold.
\end{quote}

Past (he …): \gu{તે} \gu{ટાઢમાં} \gu{થરથર્યો}, \emph{te ṭāḍhmā̃ thartharyo}.

\begin{cousinweb}
\textbf{\gu{થરથરવું}} (\emph{thartharvũ}) goes back to \emph{*thar}, \textquotedblleft{}to tremble\textquotedblright{}, a form CDIAL reconstructs and lists doubled in this verb; CDIAL lists the Gujarati verb, \textquotedblleft{}to tremble\textquotedblright{} (CDIAL 6092).
\end{cousinweb}

\begin{grammarlens}[title={What you've built}]
One more word for getting around.
\end{grammarlens}

\subsection*{Your turn}
\begin{itemize}
\raggedright
  \item \emph{Say it:} \emph{thartharvũ}
  \item \emph{Say it:} \emph{thartharvũ}, once more
  \item \emph{Say it:} \emph{ṭāḍhmā̃ bāḷko tharthartā̃ hatā̃.} — The children were shivering in the cold.
  \item \emph{Recall:} say the Gujarati for the cold, then the Gujarati for to congeal, then say \emph{thartharvũ} again
\end{itemize}

\begin{tcolorbox}[breakable,colback=teal!4,colframe=teal!35!black,title={Before you move on}]
What does \gu{થરથરવું} mean? (\textquotedblleft{}To shiver; to tremble all over\textquotedblright{}.) Say it once more.
\end{tcolorbox}
